\section{Size of the robot's movement grid in polar coordinates}
The horizontal size $X$ and vertical size $Y$ of a square of the \robot's movement
grid can be set in Cartesian coordinates by the commands
{\idt gsX(}$X${\idt )} and {\idt gsY(}$Y${\idt )}, as shown in figure \ref{fig:lgrt} below.

\begin{figure}[h]
 \includevspacePostScript[4.0cm]{LgrtExplique}{0.95}{0.95}
 \parbox{13.7cm}{\caption{The \protect\robot's movement grid defined by}\label{fig:lgrt}}
 \centering{{\idt gsX(}$X${\idt );\ gsY(}$Y${\idt )}\ or\ {\idt gsRT(}$R${\idt ,} $T${\idt )}\Large\strut}
\end{figure}

The size of a square of the grid can also be set in {\em polar
coordinates} \index{Polar coordinates}by the command
{\idt gsRT(}$R${\idt ,} $T${\idt )}, giving the length $R$ of the diagonal of this square measured
in pixels and
the angle $T$, in degrees, that this diagonal makes with a horizontal line,
as illustrated in figure \ref{fig:lgrt} above. The command {\idt gsRT} has two
parameters, which are integers or rational numbers, as in {\idt gsRT(50, 70)} or
{\idt gsRT(50.5, 22.5)}.

\begin{exemple}[Isosceles triangle]{\rm
The program below draws an isosceles triangle with two angles equal
to 70 degrees and two sides equal to 50 pixels:

{\raggedright
 %\vspace{\baselineskip}
 \vbox to 0pt{\includevspacePostScript[4cm]{TriangleIsocele}{0.95}{0.95}}%
 \input python/IsoscelesTriangle.tex
 }
\noindent $\Box$ }
\end{exemple}

\section{Rotations}
\index{Rotation}The command {\idt rot(}{\em t}{\idt )} makes the \robot\ perform a \underline{rot}ation of {\em t} degrees,
{\em t} being an integer or rational expression (for example {\idt rot(15)}, {\idt rot(-90)} or {\idt rot(22.5)}).
If the value of the expression {\em t} is positive, the robot turns on the spot
{\em counterclockwise}\index{Positive direction}\footnote{\ The counterclockwise direction is what mathematicians
call the {\em positive} (or {\em trigonometric}) {\em direction}.}. If the value of the expression {\em t} is negative, it turns on
the spot clockwise. This command does not change the position of the \robot\ but it changes
its orientation. The length of the segment drawn by the next {\idt fd()} also stays the same. To achieve this, the size of the
grid square must be changed accordingly. This is illustrated in figure \ref{fig:rtexp}
(page \pageref{fig:rtexp}), which shows a grid square in dotted lines before and after successive rotations of
15 degrees (the \robot\ starting each time from the lower left corner of the square, and the rotation being followed by an {\idt fd()} to show
that the length of the segment drawn does not change). In this figure, the
successive diagrams are arranged from top to bottom and from left to right.
\begin{figure}[htbp]
 \includevspacePostScript[17.7cm]{RTexplique}{0.95}{0.95}
 \caption{Successive rotations and moves {\idt rot(15);\ fd()}}\label{fig:rtexp}
\end{figure}

When using the command {\idt rot}, we advise you to avoid the commands {\idt etr()},
{\idt qtr()}, {\idt tetr()}, \ldots, {\idt fne()}, {\idt fse()}, \ldots, replacing them by {\idt rot} with a suitable
parameter (only the commands {\idt fn()}, {\idt fs()}, \ldots\ are useful for absolute orientations).
By using the command {\idt rot} exclusively, there is no need to take
into account the changes of coordinate system that it causes.
\begin{exemple}[Pentagon]{\rm
The program below draws a pentagon by repeating five times the drawing of a segment followed
by a rotation of $(360 / 5) = 72$ degrees:

 {\raggedright
 \vbox to 0pt{\includevspacePostScript[6cm]{Pentagone_2}{0.95}{0.95}}%
 \input python/Pentagon_2.tex
 \vspace{1.6cm}% le listing Python est plus court que la figure placee a cote
 }
\noindent $\Box$ }\end{exemple}
\begin{exemple}[Rosette of pentagons]{\rm
The rosette below is obtained by superimposing ten pentagons, each turned by 36 degrees from the previous one:

 {\raggedright
% \vspace{\baselineskip}
 \vbox to 0pt{\includevspacePostScript[6cm]{RosaceDePentagones}{0.95}{0.95}}%
 \input python/PentagonsRosette.tex
 }
\noindent $\Box$ }\end{exemple}

\section{Translations}
\index{Translation}
Whatever the orientation of the \robot, the expression {\idt valFdl()}, of type {\idt float}, has as its value the length
 (measured in pixels) of a segment that would be drawn by an {\idt fd()}.

 To draw a segment of a different length, one can use the command {\idt fdl(}$E${\idt )}
 (\underline{f}orwar\underline{d} by a \underline{l}ength), which
 takes as parameter an expression $E$ (integer or rational) whose value is the
 length of this segment
 in pixels. If this value is positive, the segment is drawn moving forward in the
 current direction of the \robot. If this value is negative, it is drawn moving backward. The \robot\
 is moved to the end of the segment and does not change its orientation. The size
 of a grid square is changed to allow this move (if the pen is up)
 or this drawing (if the pen is down). Consequently, the segment can be extended
 by a segment of the same length using an {\idt fd()} command.

 For example, the sequence of commands {\idt fdl(20);\ rot(90);\ fd();\
rot(90);\ fd();\ rot(90);\ fd()} draws a square with sides of
20 pixels.
\begin{exemple}[Polygonal Greek key]{\rm
The polygonal frieze below is obtained by replacing the drawing
of each side of the polygon by that of the Greek key (meander):

\vspace{\baselineskip}

\includevspacePostScript[6.5cm]{GrecquePolygonale}{0.95}{0.95}
\input python/PolygonalGreekKey.tex
\noindent $\Box$ }\end{exemple}
\needspace{6\baselineskip}\noindent \rule{137mm}{0.1mm}\par\nopagebreak %-----------------------------------------------------------
%%% EXOS lgrt
\begin{exercices}{Write \python\ programs using the command {\idt gsRT}
to make the following drawings:

\tiret an equilateral triangle;

\tiret the pentagon, the crosses or the spider's web of figure \ref{fig:PtCrCrTa} (page \pageref{fig:PtCrCrTa});

\tiret the geometric patterns after Horemis \cite{horemis} of figure \ref{fig:horemis} (page \pageref{fig:horemis});

\tiret the coloured tilings of figure \ref{fig:pavtrilo} (page \pageref{fig:pavtrilo}).

The very many tilings given by Grünbaum and Shephard \cite{grunbaum} are the source
of countless exercises.
\noindent $\Box$}\end{exercices}
\begin{figure}[htbp]
 \includevspacePostScript[3.1cm]{PtCrCrTa}{0.95}{0.95}
 \caption{Pentagon, crosses and spider's web}\label{fig:PtCrCrTa}
\end{figure}
\begin{figure}[htbp]
 \includevspacePostScript[6.0cm]{Horemis}{0.95}{0.95}
 \caption{Geometric patterns after Horemis}\label{fig:horemis}
\end{figure}
\begin{figure}[htbp]
 \includevspacePostScript[3.2cm]{EtoilesORCVSpi}{0.95}{0.95}
 \caption{Stars, decorated octagon and spirolateral}\label{fig:etocsp}
\end{figure}
\begin{figure}[htbp]
 \includevspacePostScript[10.2cm]{PavagesTriLos}{0.95}{0.95}
 \caption{Tilings with triangles, rhombuses and squares}\label{fig:pavtrilo}
\end{figure}
\begin{figure}[htbp]
 \includevspacePostScript[4.4cm]{RosacesPolygonales}{0.95}{0.95}
 \caption{Polygonal rosettes}\label{fig:rospol}
\end{figure}
%%% EXOS rot
\begin{exercices}{Write \python\ programs using the command {\idt rot}
to make the following drawings:

\tiret the stars, the decorated octagon or the spirolateral curve of order 3 at 135 degrees of figure
         \ref{fig:etocsp} (page \pageref{fig:etocsp}), the latter being obtained by moving forward once, then twice, then
		 three times, turning by 135 degrees each time, and repeating these moves until
		 the curve closes up;

\tiret the regular polygons of figure
         \ref{fig:polyreg} (page \pageref{fig:polyreg}), the number of sides being a parameter of the drawing
		 procedure;

\tiret the stars of figure
         \ref{fig:etoiles} (page \pageref{fig:etoiles}), the number of points being a parameter of the drawing
		 procedure;

\tiret the protractor of figure
         \ref{fig:rapport} (page \pageref{fig:rapport}), using
		 the command {\idt label(\allowbreak intToString(}{\em e}{\idt ))}, which writes the value of the integer expression {\em e}
		 to the right of the current position of the \robot\ (the irregularity of the graduations
		 is due to rounding errors in integer screen coordinates, which make it impossible to
		 draw straight lines of arbitrary slope);

\tiret the tiling of the plane with hexagons or stars of figure
         \ref{fig:pavhexeto} (page \pageref{fig:pavhexeto});

\tiret the Escher-style \cite{escher1} tiling of the disc and of the triangle of figure
         \ref{fig:pavdistri} (page \pageref{fig:pavdistri});

\tiret the tiling of the disc using Voderberg's tile \cite{voderberg} of figure
         \ref{fig:voder1} (page \pageref{fig:voder1});

\tiret the spirals using Voderberg's tile \cite{voderberg} of
         figure \ref{fig:voder2} (page \pageref{fig:voder2}).

\noindent $\Box$}\end{exercices}
\begin{figure}[htbp]
 \includevspacePostScript[7.3cm]{PolygonesReguliers}{0.95}{0.95}
 \caption{Regular polygons}\label{fig:polyreg}
\end{figure}
\begin{figure}[htbp]
 \includevspacePostScript[1.9cm]{Etoiles}{0.95}{0.95}
 \caption{Stars}\label{fig:etoiles}
\end{figure}
\begin{figure}[htbp]
 \includevspacePostScript[6.8cm]{Rapporteur}{0.95}{0.95}
 \caption{Protractor graduated in degrees}\label{fig:rapport}
\end{figure}
\begin{figure}[htbp]
 \includevspacePostScript[6.0cm]{PavageHexagonesEtoiles}{0.95}{0.95}
 \caption{Tiling of the plane with hexagons and stars}\label{fig:pavhexeto}
\end{figure}
\begin{figure}[htbp]
 \includevspacePostScript[9.7cm]{PavageDisqueTriangle}{0.95}{0.95}
 \caption{Escher-style tiling of the disc and of the triangle}\label{fig:pavdistri}
\end{figure}
\begin{figure}[htbp]
 \includevspacePostScript[6.5cm]{Voderberg1}{0.95}{0.95}
 \caption{Voderberg's tile and tiling of the disc}\label{fig:voder1}
\end{figure}
\begin{figure}[htbp]
 \includevspacePostScript[6.2cm]{Voderberg2}{0.95}{0.95}
 \caption{Voderberg spirals}\label{fig:voder2}
\end{figure}
%%% EXOS developpements de solides
\begin{exercices}{Write \python\ programs using the commands {\idt gsRT} and {\idt rot}
to draw the nets of the following solids:

\tiret the Platonic solids (regular cube, tetrahedron, octahedron, dodecahedron and
icosahedron) of figure \ref{fig:platon} (page \pageref{fig:platon});

\tiret the great and small rhombicuboctahedra and the truncated tetrahedron of figure
\ref{fig:devsol} (page \pageref{fig:devsol}).

The nets of solids given by Cundy and Rollett \cite{cundy}
are the source of many exercises.
\noindent $\Box$}\end{exercices}
\begin{figure}[htbp]
 \includevspacePostScript[5.5cm]{SolidesDePlaton}{0.95}{0.95}
 \caption{Nets of the Platonic solids}
 \label{fig:platon}
\end{figure}
\begin{figure}[htbp]
\includevspacePostScript[6.0cm]{DevelopSolides}{0.95}{0.95}
\caption{Nets of the great and small rhombicuboctahedra}\label{fig:devsol}
\centering{and of the truncated tetrahedron}
\end{figure}

%%% EXOS rot + fdl + drawToReal
\begin{exercices}{Write \python\ programs using the commands {\idt rot}
and {\idt fdl} to make the following drawings:

\tiret the polygonal rosettes of figure \ref{fig:rospol} (page \pageref{fig:rospol});

\tiret the ``jolygons'' (from the French {\em joli}, ``pretty'', and {\em polygone}) of figure \ref{fig:jgones} (page \pageref{fig:jgones}), obtained by drawing
successive segments whose length decreases by 5\,\% each time, two successive segments making
an angle of $\theta$ degrees; $\theta$ being chosen equal to 15, 60, 72, 76, 89, 92,
118, 175 and 144 degrees respectively
for the jolygons
%de la figure \ref{fig:jgones} (page \pageref{fig:jgones})
taken from left
to right and from top to bottom;

\tiret the plane curves of figure \ref{fig:courbpo} (page \pageref{fig:courbpo}),
defined by a polar equation, as follows:
 \begin{list}{$\scriptscriptstyle\bullet$}{
}
   \item the limaçon of Pascal, $\rho = a + b \cos{\theta}$ with
   $a \neq 0$, $b \neq 0$ and $\theta \in [0, 2\pi]$ (in the example, $a = 40$ and $b = 80$),
   \item the three-petalled rhodonea (rose curve), $\rho = a \cos{n \theta}$ with
    $a = 60$, $n = 3$ and $\theta \in [0, \pi]$,
   \item Moritz's curve, $\rho = a \cos{\frac{9}{10} \theta}$ with
    $a = 60$ and $\theta \in [-10 \pi, 10 \pi]$;
   \end{list}

\tiret the plane curves defined by the Cartesian equations $y = \sin{x}, x \in [-2\pi, 2\pi]$
  (figure \ref{fig:sin} below) and $y = \cos{x}, x \in [-2\pi, 2\pi]$, using the commands {\idt moveToReal}
  and {\idt drawToReal} (\underline{move to} and \underline{draw to} a point with \underline{real}-number coordinates). The command {\idt drawToReal(x, y)},
  where {\idt x} and {\idt y} are rational expressions,
  draws a segment (if the pen is down) from the current position of the
  \robot\ to the point with Cartesian coordinates $(x,\ y)$ and moves the \robot\ to
  this new position $(x,\ y)$. The command {\idt moveToReal(x, y)} does the
  same without drawing. The functions {\idt sin} and {\idt cos} and the constant {\idt pi}
  are explained in the addendum of chapter \ref{ch:rec}.
   Many other examples of curve drawing are given in \cite{LT}.

\noindent $\Box$}\end{exercices}
\begin{figure}[htbp]
 \includevspacePostScript[4.7cm]{Sinus}{0.95}{0.95}
 \caption{Cartesian curve $y = \sin{x}, x \in [-2\pi, 2\pi]$}\label{fig:sin}
\end{figure}
\begin{figure}[htbp]
 \includevspacePostScript[12.0cm]{Jolygones}{0.95}{0.95}
 \caption{Jolygons at 15, 60, 72, 76, 89, 92, 118, 175 and 144 degrees}\label{fig:jgones}
\end{figure}
\begin{figure}[htbp]
 \includevspacePostScript[5.3cm]{CourbesPolaires}{0.95}{0.95}
 \caption{Limaçon of Pascal, three-petalled rhodonea, Moritz's curve}\label{fig:courbpo}
\end{figure}
%%% CORRIGES lgrt
\begin{corrige}{
We give the program that draws the spider's web of figure \ref{fig:PtCrCrTa}
(page \pageref{fig:PtCrCrTa}), the other programs being in the folder
{\idt book/\allowbreak 20\_polar\_grid}:
\input python/SpiderWeb.tex
\noindent $\Box$}\end{corrige}
%%% CORRIGES rot
\begin{corrige}{
We give the procedure that draws a regular polygon with $n$ sides (figure \ref{fig:polyreg},
page \pageref{fig:polyreg}), the other programs being in the folder
{\idt book/\allowbreak 20\_polar\_grid}:
\input python/RegularPolygon.tex
\noindent $\Box$}\end{corrige}
%%% CORRIGES developpements de solides
\begin{corrige}{
We give the program that draws the net of the truncated tetrahedron, the other
programs being in the folder {\idt book/\allowbreak 20\_polar\_grid}:
\input python/TruncatedTetrahedron.tex
\noindent $\Box$}\end{corrige}
%%% CORRIGES rot + fdl
\begin{corrige}{
We give the program that draws the jolygon at 175 degrees; the other jolygons are obtained by changing the
value of the constant {\idt angle}. The other programs are in the folder
{\idt book/\allowbreak 20\_polar\_grid}.
\input python/Jolygon175.tex
\noindent $\Box$}\end{corrige}

\clearpage

%%%%%%%%%%%%%%%%%%%%%%%%%%%%%%%%%%%%%%%%%%%%%%%%%%%%%%%%%%%%%%%%%%%%%%%%%%%
\noindent \framebox[\textwidth]{
\begin{minipage}{\myparboxwidth}
\setlength{\parindent}{2em}
\small \sf
\vspace{\abovedisplayskip}
\section*{Addenda}
The command {\idt gsRT(r, t)} sets the length {\idt r} of the diagonal of a grid
square (in pixels) and the angle {\idt t}, in degrees, that this diagonal makes
with a horizontal line. It is equivalent to:
\begin{center}
{\idt gsX(r * cos(t / 180 * pi))}\\
{\idt gsY(r * sin(t / 180 * pi))}
\end{center}
but avoids using trigonometric functions.
The angle {\idt t}, between 0 and 90, defines the north-east orientation, which coincides with east when {\idt t} is {\idt 0}
and with north when {\idt t} is {\idt 90}.

The command {\idt rot(t)} makes the \robot\ turn on the spot through an angle of {\idt t} degrees to the left,
in the positive (trigonometric) direction, that is counterclockwise, if {\idt t} $> 0$, and to the right, clockwise,
if {\idt t} $ < 0$. The command {\idt fdl(l)} makes the \robot\ move straight forward (if {\idt l} $> 0$)
or backward (if {\idt l} $< 0$) by {\idt l} times the length of the side of a pixel
of the screen.

These commands {\idt rot} and {\idt fdl} are directly inspired by the ``turtle'' graphics of
{\sc logo}. They change the coordinate system and the orientation of the \robot\ so as to remain compatible
with the use of Cartesian or polar coordinates.

The {\sc logo} language was designed by Seymour Papert at the end of the 1970s \cite{logo}.
It has many similarities with {\sc lisp}
\cite{lisp}, in particular the interpretation of data representing instructions,
dynamic scoping of identifiers and the list structure.
Interpreting data as programs entails some subtleties in
writing programs (such as the use of the {\idt quote} function, \cite{wilson}) which are not
within the reach of all beginners in computer science.
Dynamic scoping of identifiers \cite{wilson}
has largely been abandoned, and the list structure
is almost always supplemented by arrays in modern {\sc lisp}s such as
{\sc common lisp}, which in this follow the languages descended from {\sc algol 60}
\cite{algol60}, such as \pascal\ or {\sc ada}. The \python\ language also offers
both lists and arrays, does not interpret its data as
instructions (unless explicitly asked to) and uses static scoping of identifiers.

\vspace{\abovedisplayskip}
\end{minipage}
}
%%%%%%%%%%%%%%%%%%%%%%%%%%%%%%%%%%%%%%%%%%%%%%%%%%%%%%%%%%%%%%%%%%%%%%%%%%%%
\newpage                 % Si on termine le chapitre par une page blanche
\thispagestyle{empty}
